\section{Bài toán, phạm vi và yêu cầu}
\subsection{Bối cảnh và mục tiêu}
Dự án dùng bộ dữ liệu giao dịch thương mại điện tử OList để tổ chức một kho dữ liệu phục vụ quản lý kinh doanh. Hệ thống gồm ba lớp: (1) kho dữ liệu đa chiều nạp từ CSV qua staging; (2) phân tích OLAP qua API Flask và giao diện React; (3) một DSS dự báo nhu cầu tháng tới và đề xuất lượng nhập hàng.

Các nhu cầu nghiệp vụ gồm xác định danh mục, thị trường và người bán đóng góp lớn, theo dõi doanh số theo thời gian, đánh giá khách mua lại và hỗ trợ quyết định nhập hàng. Phần đánh giá ban đầu về phân khúc khách hàng (RFM, K-Means) đã được thử trên dữ liệu thật nhưng cho tín hiệu yếu; lý do chọn hướng dự báo nhu cầu được trình bày ở chương~\ref{sec:dss}.

\subsection{Đối chiếu yêu cầu và kết quả}
\begingroup\small\renewcommand{\arraystretch}{1.18}
\begin{longtable}{>{\raggedright\arraybackslash}p{0.22\linewidth}>{\raggedright\arraybackslash}p{0.40\linewidth}>{\raggedright\arraybackslash}p{0.30\linewidth}}
\caption{Đối chiếu yêu cầu và kết quả}\label{tab:requirements}\\
\toprule Yêu cầu & Kết quả & Chương\\\midrule\endfirsthead
\toprule Yêu cầu & Kết quả & Chương\\\midrule\endhead
\bottomrule\endfoot
Nghiệp vụ & Chín câu hỏi phân tích (bảng~\ref{tab:bq}) và một câu hỏi quyết định & 1\\
DW đa chiều, ETL & Hai fact, bốn dimension; raw $\rightarrow$ staging $\rightarrow$ DW; nghiệm thu F1--F7 & 3, 4\\
Backend và OLAP & Flask; ba nhóm API Thời gian, Khu vực/Người bán, Khách hàng & 5\\
Frontend & Sáu màn: Tổng quan, Thời gian, Khu vực, Người bán, Khách hàng, Quyết định & 5, 6\\
Thành phần DSS & Dự báo nhu cầu theo tháng, lượng nhập đề xuất, đường đánh đổi thiếu/dư hàng & 6\\
Đánh giá kết quả & Đối soát tổng tiền, KPI, backtest DSS & 7\\
\bottomrule
\end{longtable}\endgroup

\subsection{Câu hỏi phân tích}
\begin{table}[H]\centering\small\renewcommand{\arraystretch}{1.15}
\begin{tabularx}{\linewidth}{lYY}
\toprule Mã & Câu hỏi nghiệp vụ & Thước đo và nơi xử lý\\\midrule
BQ1 & Tiền hàng theo tháng, quý, năm thay đổi thế nào? & Doanh thu hàng, MoM/QoQ/YoY; màn Thời gian.\\
BQ2 & Danh mục nào đóng góp tiền hàng lớn nhất? & Doanh thu theo danh mục; màn Khu vực.\\
BQ3 & Người bán nào tạo nhiều tiền hàng? & Doanh thu theo seller; màn Người bán.\\
BQ4 & Bang, thành phố nào mua nhiều? & Doanh thu và số đơn theo địa lý khách; màn Khu vực.\\
BQ5 & Cước vận chuyển chiếm tỷ lệ bao nhiêu so với giá? & Cước/Doanh thu hàng; Tổng quan và Khu vực.\\
BQ6 & Giá trị thanh toán trung bình mỗi đơn là bao nhiêu? & AOV trên đơn có dữ liệu thanh toán.\\
BQ7 & Bao nhiêu khách mua lại? & Số đơn delivered theo khách duy nhất; màn Khách hàng.\\
BQ8 & Khách mua và chi tiêu thế nào theo bang, năm, trạng thái đơn? & Drill-down, dice, ma trận bang $\times$ năm.\\
BQ9 & Các trạng thái đơn phân bố ra sao? & Đếm mọi trạng thái; tiền chỉ tính đơn delivered.\\
BQ10 & Tháng tới bán được bao nhiêu và nên nhập bao nhiêu? & DSS, màn Quyết định.\\
\bottomrule\end{tabularx}
\caption{Câu hỏi phân tích}\label{tab:bq}
\end{table}

\subsection{Phạm vi dữ liệu và khái niệm doanh số}
Kho lưu mọi trạng thái đơn. Các KPI bán hàng và khách hàng mặc định lọc \code{order_status = 'delivered'}; thống kê trạng thái dùng toàn bộ đơn. Tiền hàng (\code{price}), cước (\code{freight_value}) và thanh toán (\code{total_payment_value}) là ba thước đo riêng. ``Doanh thu hàng'' là tổng \code{price}, không đồng nhất với tiền thanh toán hay lợi nhuận; dữ liệu không có giá vốn.

Loại thanh toán (\code{payment_type}), số kỳ trả góp, đánh giá (review) và geolocation chỉ có ở CSV, không được đưa vào staging có kiểu hay DW. Vì vậy câu hỏi về phương thức thanh toán và các phân tích dựa trên review hay thời gian giao hàng nằm ngoài phạm vi hệ thống.
